\subsection{The complete two-connector coefficient certificate}\label{proof:Q5-C037-coefficient}
We prove the pointwise premise \eqref{eq:L37:C2} for the same
two-connector message, with the original row constraints retained.
Its vector $R=r+t_3=K(g+s)$ has $\E(g+s)=0$; centering an original
row gives
\[
|R|\le\sqrt{pq(d-2\delta_2+\delta_3)}
\le\sqrt{pqd[1-(51/32)v]},\qquad |r|\le\sqrt{pqd}.
\]
Put $x=\sqrt d$, $t=\sqrt{d[1-(51/32)v]}$. Then
\[
a=1-(35/51)x^2-(16/51)t^2,\qquad
m=1-(18/17)x^2-(16/17)t^2,
\]
where $m=\mu_2$. The actual parameters satisfy $(19/32)x\le t\le x$
and $t^2\le d-(51/32)d^2\le8/51<4/25$.
It suffices to work on the full closed rectangle
\begin{equation}\tag{L37.8.3}\label{eq:L37:rectangle}
0\le x\le1/\sqrt3,\qquad0\le t\le2/5,
\end{equation}
where $a,m>0$. Since $\sqrt{pq}\le1$, the actual coefficient is at most
\[
(p+x)\{[a(2-p)+t]^2+2m[a(2-p)+t]+3m^2\}.
\]
Writing $y=p+x\in[x,2+x]$ and $B_0=a(2+x)+t$ makes this
$H(y)=y[(B_0-ay)^2+2m(B_0-ay)+3m^2]$.
At $y_0=7/5$, define
\[
\begin{gathered}
z_0=a(3/5+x)+t,\quad u_0=z_0+m,\quad H_0=(7/5)(u_0^2+2m^2),\\
D=H'(7/5)=u_0^2+2m^2-(14/5)au_0.
\end{gathered}
\]
The proof will establish $|D|<3/10$ on the entire rectangle, and the
additional independent estimates needed to combine that derivative bound.

\proofpart{The exact Taylor square for every allowed value of the weight}
For $h_0=y-7/5$, the cubic expansion is exactly
\[
H(y)=H_0+Dh_0-a[2u_0-ay]h_0^2.
\]
Since $y\le2+x$, the square coefficient is at least
$K_0=a[a(x-4/5)+2t+2m]$.
Let $b=16/51$, $A_0=1-35x^2/51$, $M_0=1-18x^2/17$,
$L=x+26/5$, $B_1=A_0(x-4/5)+2M_0$. Then
\[
K_0=(A_0-bt^2)(B_1+2t-bLt^2),\quad
\partial_t^2K_0=-2b(A_0L+B_1)-12bt+12b^2Lt^2<0.
\]
Indeed $A_0>3/4$, $L\le29/5$, $B_1\le2$, $A_0L+B_1>4$;
with $b<1/3$, $t\le2/5$ the last derivative is less than
$-8b+(1392/125)b^2<0$.
Dropping positive terms from the endpoint difference and using
$A_0L+B_1\le39/5$ gives
$K_0(x,\allowbreak 2/5)-K_0(x,\allowbreak 0)>3/5-52/125-16/375>0$.
Thus the smaller $t$ endpoint is zero. There the exact polynomial is
\[
k(x)=6/5+x-(122/51)x^2-(70/51)x^3
+(2800/2601)x^4+(1225/2601)x^5.
\]
Its second derivative is convex and negative at both endpoints:
$k''(0)=-244/51$ and
$k''(1/\sqrt3)=-1244/2601-39760/(7803\sqrt3)<0$.
So $k$ is concave. Its endpoint minimum exceeds $13/16$:
$k(0)=6/5$, and using $1/\sqrt3>1/2$ gives
\[
k(1/\sqrt3)>(118/153)[59/153+518/765]=31978/39015>13/16.
\]
Consequently $K_0>13/16$ on the full rectangle, without requiring
$H''\le0$ for all $y$, and
\begin{equation}\tag{L37.8.6}\label{eq:L37:Taylor}
H(y)\le H_0+(4/13)D^2\le H_0+D^2/3.
\end{equation}

\proofpart{The complete anchor-value certificate}
Let $F(x,t)=H_0+8x^2$, $b=16/51$, $q_0=x+18/5$,
\[
U_0=8/5+x-(25/17)x^2-(35/51)x^3,\qquad M_0=1-(18/17)x^2.
\]
Then $u_0=U_0+t-bq_0t^2$, $m=M_0-3bt^2$.
Concavity of $U_0$ and its endpoint values give $U_0\ge3/2$;
also $M_0\ge11/17$ and $18/5\le q_0\le21/5$.
The exact second derivative is
\[
\frac{F_{tt}}{14/5}=1-2bq_0U_0-12bM_0-6bq_0t
+b^2(6q_0^2+108)t^2.
\]
The constant part is at most $1-288/85-704/289<-19/4$.
The remaining part is at most $6/25$: the convex polynomial
$R(q)=-6bq+(2/5)b^2(6q^2+108)$ has endpoint values
$174528/325125$ and $166752/325125$, both less than $3/5$,
and is multiplied by $t\le2/5$. Hence $F_{tt}<-12$.
At $t=1/4$,
\[
\begin{gathered}
F_t(x,1/4)=(14/5)N(x),\\
N=-643/5780+(128/867)x+(527/2601)x^2
-(59/867)x^3+(280/2601)x^4.
\end{gathered}
\]
The numerator of $N'$ is $384+1054x-531x^2+1120x^3>0$, using
$x^2\le1/3$. Thus $F_t$ increases in $x$. Its left endpoint is
$-4501/14450>-1/3$; at the right endpoint, $1/\sqrt3<3/5$ yields
$N<-1/10+7/100+3/250+3/40=57/1000$ and hence $F_t<1/3$.
So $|F_t(x,1/4)|<1/3$, giving
$F(x,t)<F(x,1/4)+1/216$ by strong concavity.

Set $F_0(x)=F(x,1/4)$. Its complete derivative is
\[
F_0''=-92647/13005-(12929/289)x+(222796/4335)x^2
+(49000/867)x^3+(17150/867)x^4.
\]
All coefficients of $F_0^{(4)}$ are positive, so $F_0''$ is convex.
The left endpoint is below $-2$; the right is
\[
F_0''(1/\sqrt3)=476197/39015-67361/(2601\sqrt3)
<49/4-(129/5)(4/7)=-349/140<-2.
\]
At $2/5$,
$F_0=115038439/14450000<637/80$ and
$F_0'=-24517/541875\in(-1/20,0)$.
Thus $F_0<637/80+1/1600$, and the separate strong concavity bounds give
\[
F<637/80+1/1600+1/216=255/32-43/43200.
\]
For the stronger combined bound, put $n(x)=F_t(x,1/4)$.
The same numerator for $N'$ gives $0<n'(x)<7/5$: discarding its
nonpositive quadratic term and using $x<3/5$ bounds it by
\[
\frac{14}{5}\frac{384+1054(3/5)+1120(3/5)^3}{2601}
=440412/325125<7/5.
\]
Also $n(2/5)=-12971/216750\in(-1/16,0)$.
For $h=x-2/5$, $z=t-1/4$, $s=F_0'(2/5)$ and $n_0=n(2/5)$,
the two strong concavity bounds and this Lipschitz bound give
\[
\begin{aligned}
F(x,t)&\le F_0(2/5)+sh+n_0z-h^2-6z^2+(7/5)|hz|\\
&<637/80+sh+n_0z-h^2/2-5z^2.
\end{aligned}
\]
Here $(7/5)|hz|\le h^2/2+(49/50)z^2$.
Completing both squares yields
\begin{equation}\tag{L37.8.9}\label{eq:L37:anchor}
H_0+8d<637/80+1/800+1/5120=203877/25600.
\end{equation}

\proofpart{The derivative certificate on the entire rectangle}
Retain the full polynomial certificate for the same
$D=u_0^2+2m^2-(14/5)au_0$. Write
$D(x,t)=\sum_{j=0}^4D_j(x)t^j$. Direct expansion gives
\begin{align*}
D_0={}&2/25+(2/5)x-(191/255)x^2-(22/17)x^3\\
&+(179/867)x^4+(1820/2601)x^5+(1225/2601)x^6,\\
D_1={}&2/5+2x-(52/51)x^2-(70/51)x^3,\\
D_2={}&-2309/1275-(128/85)x+(13952/4335)x^2
+(3296/2601)x^3+(1120/2601)x^4,\\
D_3={}&-2b(x+11/5),\\
D_4={}&b^2\mathcal C(x),\quad\mathcal C=x^2+(22/5)x+522/25.
\end{align*}
Indeed, with $A_0=1-35x^2/51$, $M_0=1-18x^2/17$ one has
$u_0=U_0+t-b(x+18/5)t^2$, giving all five coefficients.
These are ordinary polynomials, including the rectangle boundary.

First,
\[
D_{tt}=2D_2-12b(x+11/5)t+12b^2\mathcal C(x)t^2.
\]
It is convex in each coordinate separately, since $D_2''>0$ and
$\mathcal C''=2>0$. Its maximum is therefore at a corner.
Put $X=1/\sqrt3$. The exact value
\[
D_2(X)=-403991/585225-718624/(663255\sqrt3)<-9/7
\]
follows because the first term is below $-2/3$, the coefficient in
the second exceeds $13/12$ (the integer difference is $1173$), and
$X>4/7$. At $(0,0)$, $D_{tt}=-4618/1275<-18/5$; at $(X,0)$ it is
less than $-18/7$. At $(0,2/5)$ its three terms are less than
$-18/5,-33/10,4$, summing to $-29/10$.
At $(X,2/5)$ the first is below $-18/7$; since $X>1/2$ the negative
term has absolute value greater than
$(24/5)(16/51)(27/10)=1728/425>81/20$.
Since $b^2<1/10$, $X<3/5$, $X^2=1/3$,
$\mathcal C(X)<1/3+66/25+522/25=1789/75$,
and the positive term is below
$(48/25)(1/10)(1789/75)=85872/18750<23/5$.
Thus the final corner is below
$-18/7-81/20+23/5=-283/140<-2$, proving $D_{tt}<-2$ throughout.

Second, the full other derivative is
\begin{align*}
D_{xx}={}&-382/255-(132/17)x+(716/289)x^2
+(36400/2601)x^3+(36750/2601)x^4\\
&+[-104/51-(140/17)x]t\\
&+[27904/4335+(19776/2601)x+(13440/2601)x^2]t^2
+(512/2601)t^4.
\end{align*}
It too is separately convex: in $t$ directly, and in $x$ because its
quadratic and higher coefficients are nonnegative, with positive
quadratic coefficient. At $(0,0)$ it equals $-382/255<-4/5$.
At $(X,0)$,
\[
D_{xx}=105072/117045-24188/(7803\sqrt3)
<9/10-3(4/7)=-57/70<-4/5.
\]
At $(0,2/5)$ it is
\[
-118/51+111616/108375+8192/1625625.
\]
The three terms are below $-23/10$, $21/20$ and $1/100$, and their
sum is less than $-4/5$.
For the last corner compare to $(X,0)$. Its linear $t$ term is less
than $-(2/5)(104/51+70/17)=-628/255<-12/5$.
The bracket in the quadratic $t$ term is less than
$13/2+24/5+7/4=261/20$, using $X<3/5$ and $X^2=1/3$;
after multiplication by $4/25$ it is below $21/10$.
The quartic term is below $1/100$.
The increment is therefore below $-12/5+21/10+1/100=-29/100<0$.
All corners satisfy the desired bound, proving $D_{xx}<-4/5$.

For an upper bound on $D$, take $x_0=1/4$, $t_0=1/5$.
There $a=6423/6800$, $m=3047/3400$, $u_0=258271/136000$ and
\[
\begin{gathered}
D(x_0,t_0)=3516398793/18496000000<191/1000,\\
D_t(x_0,t_0)=-7468703/86700000,\quad |D_t|<1/10.
\end{gathered}
\]
The $t$ concavity bound gives
$D(1/4,t)<191/1000+1/400=387/2000<1/5$ on the full vertical line.
On that line the exact coefficient derivatives are
\[
\begin{gathered}
D_1'=503/408\in(6/5,5/4),\quad D_2'=4784/13005\in(0,2/5),\\
D_3'=-32/51<-3/5,\quad D_4'=6272/13005\in(0,1/2),\\
D_0'=13/510-33/136+179/13872+2275/166464+1225/443904
\in(-1/5,-9/50).
\end{gathered}
\]
For the last interval the lower bound uses $13/510>1/40$,
$33/136<1/4$ and the next two positive terms each greater than $1/80$;
the upper uses $13/510<13/500$, $33/136>6/25$, and the remaining
positive terms less than $13/1000,7/500,3/1000$.
Furthermore
\[
\partial_tD_x(1/4,t)=D_1'+2D_2't+3D_3't^2+4D_4't^3
>6/5-(96/51)(4/25)=382/425>0.
\]
Thus its lower endpoint exceeds $-1/5$, and its upper endpoint is below
\[
-9/50+(5/4)(2/5)+(2/5)(4/25)-(3/5)(8/125)+(1/2)(16/625)
=224/625<3/8.
\]
Hence $|D_x(1/4,t)|<3/8$. The $x$ strong concavity bound yields
\[
D(x,t)<1/5+(3/8)^2/(8/5)=1/5+45/512=737/2560<3/10.
\]
Maximizing that quadratic on the whole real line merely enlarges the
allowed interval; it requires no assumption on the maximizing point.

For the lower bound, the established separate concavity puts the minimum
at one of the four corners, without asserting joint concavity.
At $(0,0)$, $D=2/25>0$.
At $(0,2/5)$,
$a=1211/1275<19/20$, $m=361/425>21/25$,
$u_0=3866/2125\in(9/5,11/6)$, so
\[
D>81/25+2(441/625)-(14/5)(19/20)(11/6)
=-1691/7500>-3/10.
\]
At $(X,0)$, $a=118/153<31/40$, $m=11/17>16/25$;
$4/7<X<7/12$ gives
$8303/5355<u_0<7159/4590$ and thus $31/20<u_0<25/16$.
Therefore
\[
D>(31/20)^2+2(16/25)^2-(14/5)(31/40)(25/16)
=-6757/40000>-3/10.
\]
At $(X,2/5)$, $a=2758/3825<73/100$, $m=211/425>49/100$;
the same bounds on $X$ give
$233093/133875<u_0<401558/229500$ and $87/50<u_0<7/4$.
Thus
\[
D>(87/50)^2+2(49/100)^2-(14/5)(73/100)(7/4)
=-173/2500>-3/10.
\]
This proves $-3/10<D<3/10$ on every point of the full closed rectangle.
Finally \eqref{eq:L37:Taylor} and \eqref{eq:L37:anchor} give
\[
H(y)+8d<203877/25600+(4/13)(3/10)^2
=1023/128-183/332800.
\]
In particular the required coefficient $C_2(d)=1023/128-8d$ in
\eqref{eq:L37:C2} is valid, completing the two-connector premise.
